\section{模型假设}

\begin{enumerate}[label=\arabic*.]
\item \textbf{转写顺序一致性。} 假设官方转写与视频语音中的词序保持一致，话语的先后关系在转写中完整保留。该假设使词级强制对齐可以在固定词序下搜索单调路径；无法构造字符目标的词仍按同一顺序单独记录，不参与路径搜索。

\item \textbf{词级时间共享性。} 假设同一原词对应的子词共享父词的时间区间，音视频特征在该区间内统一聚合，子词之间不再细分时间。依据是词为发音的自然单位：原词给出比子词更稳定的时间锚点，子词只改变序列长度，不改变时间归属，接口位置可经父词索引回溯到同一个时间区间。

\item \textbf{局部稳定性。} 假设音频与视觉帧特征在其短时间支撑区间内近似稳定，因此可以采用时间交叠长度作为词级聚合权重，把锚点区间上的连续积分退化为按交叠时长的加权求和。

\item \textbf{缺失不改变标签。} 假设问题二中由人工遮挡构造的局部缺失只改变模型可获取的观测信息，不改变样本对应的真实情感极性与强度；遮挡后的监督目标仍取原样本标签，以评价输入信息减少时的预测性能。

\item \textbf{视觉线索由人脸区域承载。} 假设口型、眉眼形态与面部运动足以表征样本中与情感判断相关的视觉信息，因此视觉分支只取人脸检测框内的人脸几何与关键点运动，画面中的场景、字幕与肢体动作不进入模型输入。

\item \textbf{缺失机制与标签独立。} 假设某一位置是否缺失不依赖于该样本的情感极性与强度：附件2的自然缺失来自采集与检测过程，附件3的缺失按赛题规则构造，缺失状态本身不携带情感信息。模型因此只从观测证据学习条件分布，两种来源的缺失可以在同一掩码框架下处理。
\end{enumerate}
